\documentclass[11pt]{amsart}
\usepackage[T1]{fontenc}
\usepackage{lmodern,amsmath,amssymb,amsthm,microtype}
\usepackage[margin=1.15in]{geometry}
\usepackage[colorlinks=true,linkcolor=blue,citecolor=blue,urlcolor=blue]{hyperref}
% Repository locator for the public (non-anonymous) build of
% non_mf_groups_exist.tex.  Kept in a separate file so that the anonymous
% build has an anonymous source archive and not merely an anonymous PDF:
% exclude this file from a double-blind submission and the manuscript still
% compiles, printing "supplied with this submission as supplementary
% material" in place of the link.
%
% Loaded only when \anonymousfalse is set in the preamble.
\newcommand{\coderepository}{https://github.com/SauersML/group-approximation}
\newcommand{\coderepositoryname}{github.com/SauersML/group-approximation}

\hypersetup{pdftitle={Infinite simple Kazhdan groups that are sofic},pdfauthor={\paperauthorname}}
\newtheorem{theorem}{Theorem}
\newtheorem{proposition}[theorem]{Proposition}
\newtheorem{corollary}[theorem]{Corollary}
\newtheorem{question}{Question}
\newcommand{\F}{\mathbb F}
\newcommand{\Z}{\mathbb Z}
\newcommand{\LC}{\operatorname{LC}}
\newcommand{\EL}{\operatorname{EL}}
\newcommand{\GL}{\operatorname{GL}}
\newcommand{\PSL}{\operatorname{PSL}}
\newcommand{\Sym}{\operatorname{Sym}}
\newcommand{\diag}{\operatorname{diag}}
\newcommand{\doi}[1]{\href{https://doi.org/#1}{\nolinkurl{doi:#1}}}
\raggedbottom
\emergencystretch=1.5em

\begin{document}
\title{Infinite simple Kazhdan groups that are sofic}
\author{\paperauthorname}
\date{}
\subjclass[2020]{Primary 20E32; Secondary 20E26, 20F10, 22D10, 37B10, 16S35}

\begin{abstract}
Let a finitely generated group $\Gamma$ act freely on an infinite minimal
subshift $X$ whose patterns on every finite window are those of a finite
subshift. Then for $n\ge3$ the group $\EL_n(\LC(X,\F_2)\rtimes\Gamma)$ is
infinite, finitely generated and simple, with property~\textup{(T)}. It
is locally embeddable into finite groups, so it is sofic and
hyperlinear. Such subshifts exist exactly over residually finite groups,
and for $\Gamma=\Z$ every infinite minimal subshift qualifies. This
answers the question of Brown and Ozawa whether an infinite simple
Kazhdan group can be hyperlinear, and Pestov's sofic version of it.
Every finitely generated residually finite group is a subgroup of one of
these groups, and every Turing degree occurs as the word-problem degree
of one of them.
\end{abstract}
\maketitle

Can an infinite simple Kazhdan group be hyperlinear? Brown asked this
in its von Neumann algebra form in 2001~\cite[\S11, Question~7]{Brown},
and Ozawa in the hyperlinear form in 2003~\cite[p.~527]{Ozawa}.
Pestov's Open question~9.1 adds the sofic version~\cite{Pestov}.
The following construction answers all three forms positively.

Let $\Gamma$ be a finitely generated group and $A$ a finite set. A
subshift is a closed subset $X\subseteq A^\Gamma$ invariant under the
action $(\gamma x)(h)=x(h\gamma)$. For finite $F\subseteq\Gamma$ put
$L_F(X)=\{x|_F:x\in X\}$. We say that $X$ has \emph{finite models} if for
every finite $F\subseteq\Gamma$ some finite subshift $Y$ has
$L_F(Y)=L_F(X)$.

\begin{theorem}\label{thm:main}
Let $\Gamma$ act freely on an infinite minimal subshift $X\subseteq A^\Gamma$
with finite models, and let $\LC(X,\F_2)$ be the ring of locally constant
functions $X\to\F_2$. Then for every $n\ge3$ the group
\[
  \EL_n\bigl(\LC(X,\F_2)\rtimes\Gamma\bigr)
\]
is an infinite, finitely generated, simple group with Kazhdan's
property~\textup{(T)}. It is locally embeddable into finite groups
\textup{(LEF)}, so it is sofic and hyperlinear.
\end{theorem}

\begin{proposition}\label{prop:models}
\begin{enumerate}
\item[\textup{(a)}] Every infinite minimal subshift $X\subseteq A^{\Z}$ is
free and has finite models.
\item[\textup{(b)}] If $\Gamma$ acts freely on a nonempty subshift with
finite models, then $\Gamma$ is residually finite.
\item[\textup{(c)}] Every infinite finitely generated residually finite
group acts freely on an infinite minimal subshift with finite models.
\end{enumerate}
\end{proposition}

For an infinite minimal subshift $X\subseteq A^\Z$ we write
$G_X=\EL_3(\LC(X,\F_2)\rtimes\Z)$. By Theorem~\ref{thm:main} and
Proposition~\ref{prop:models}(a), every $G_X$ is an infinite simple
Kazhdan group that is sofic and hyperlinear.

\begin{corollary}\label{cor:rf}
Every finitely generated residually finite group is a subgroup of an
infinite finitely generated simple Kazhdan LEF group.
\end{corollary}

Property~\textup{(T)} follows from the theorem of Ershov and
Jaikin-Zapirain~\cite{EJZ}. For $\Gamma=\Z$ the finite models are
periodic sequences with the same short words as $X$, as in Grigorchuk
and Medynets~\cite{GM}. Kerr and Nowak studied residually finite actions
and their crossed products~\cite{KerrNowak}. The new step is simplicity.
A nontrivial normal subgroup contains a nontrivial commutator lying in a
finite simple group $\GL_d(\F_2)$ acting on a clopen tower, so it
contains this group and with it an elementary matrix. For derived
topological full groups, Matui showed in the same way that a nontrivial
normal subgroup meets a simple union of alternating groups on
towers~\cite[Lemma~3.4 and Theorem~4.9]{Matui}. Thom constructed a
finitely generated Kazhdan LEF group that is not residually
finite~\cite{Thom}, but his example is not simple. Without
property~\textup{(T)}, Corollary~\ref{cor:rf} is due to Kionke and
Schesler~\cite[Theorem~1.2]{KionkeSchesler}.

\section{Proof of Theorem~\ref{thm:main}}

\subsection*{The ring and property (T)}
Let $S$ be a finite symmetric generating set of $\Gamma$, with word length
$|\gamma|$ and balls $B_w=\{\gamma:|\gamma|\le w\}$. Write $e_U$ for the
indicator of a clopen set $U\subseteq X$, and
\[
 R=\LC(X,\F_2)\rtimes\Gamma=\bigoplus_{\gamma\in\Gamma}\LC(X,\F_2)\,u_\gamma,
 \qquad u_\gamma fu_\gamma^{-1}=f\circ\gamma^{-1},\quad
 u_\gamma u_\delta=u_{\gamma\delta}.
\]
So every element of $R$ is uniquely a finite sum
$\sum_\gamma f_\gamma u_\gamma$, and $u_\gamma e_Uu_\gamma^{-1}=e_{\gamma U}$.
For $\Gamma=\Z$ and $u=u_1$ the action is $(Tx)_t=x_{t+1}$, and
$ufu^{-1}=f\circ T^{-1}$.

The ring $R$ is generated by the $u_t$, $t\in S$, and the letter
indicators $e_a=e_{\{x:x(e)=a\}}$, since products of the translates
$u_\gamma e_au_\gamma^{-1}$ are the cylinder indicators, which span
$\LC(X,\F_2)$. Let $G=\EL_n(R)$ be the subgroup of $\GL_n(R)$ generated by
the elementary matrices $e_{ij}(r)=I_n+rE_{ij}$, $i\ne j$, $r\in R$, where
$E_{ij}$ is the usual matrix unit. With $[g,h]=ghg^{-1}h^{-1}$,
\begin{equation}\label{eq:elementary}
 \begin{aligned}
 e_{ij}(r+s)&=e_{ij}(r)e_{ij}(s),\qquad i\ne j,\\
 [e_{ik}(r),e_{kj}(s)]&=e_{ij}(rs),\qquad i,j,k\text{ distinct}.
 \end{aligned}
\end{equation}
So the matrices $e_{ij}(s)$ with $s\in\{u_t:t\in S\}\cup\{e_a:a\in A\}$
generate $G$. By Ershov and Jaikin-Zapirain~\cite[Theorem~1.1]{EJZ},
$\EL_n$ of a finitely generated associative unital ring has
property~\textup{(T)} for $n\ge3$. So $G$ has property~\textup{(T)}. It is
infinite because $e_{12}(\LC(X,\F_2))$ is infinite.

\subsection*{Finite models}
A group is LEF~\cite{VershikGordon} if every finite subset embeds
injectively into a finite group, preserving all products that stay in
that subset.

Let $\gamma\ne e$. The clopen sets $\{x:x(h)\ne x(h\gamma)\}$, $h\in\Gamma$,
cover $X$ because the action is free. By compactness some finite
$F_\gamma\subseteq\Gamma$ has the property that every $x\in X$ satisfies
$x(h)\ne x(h\gamma)$ for some $h\in F_\gamma$. So no point of a finite
subshift with the patterns of $X$ on $F_\gamma\cup F_\gamma\gamma$ is fixed
by $\gamma$.

Let $Y$ be a finite subshift and $F\subseteq\Gamma$ finite with
$L_F(Y)=L_F(X)$. On $\F_2^Y$ let $P_\gamma\delta_y=\delta_{\gamma y}$. For
$f\in\LC(X,\F_2)$ depending only on the coordinates in $F$ let
$D(f)\delta_y=f(x)\delta_y$, where $x\in X$ agrees with $y$ on $F$, and put
$\varphi_Y(\sum_\gamma f_\gamma u_\gamma)=\sum_\gamma D(f_\gamma)P_\gamma$.
If $f$ depends only on the coordinates in $W$ and
$W\cup W\gamma^{-1}\subseteq F$, then $P_\gamma D(f)P_\gamma^{-1}=D(f\circ\gamma^{-1})$.
So for $r,s\in R$ the identities $\varphi_Y(r+s)=\varphi_Y(r)+\varphi_Y(s)$
and $\varphi_Y(rs)=\varphi_Y(r)\varphi_Y(s)$ hold once $F$ is large. Let
$r=\sum_\gamma f_\gamma u_\gamma\ne0$ have support $E$, and let $F$ also
contain $F_g\cup F_gg$ for all $g\in E^{-1}E\setminus\{e\}$. Then for each
$y\in Y$ the points $\gamma y$, $\gamma\in E$, are distinct. If
$f_{\gamma_0}\ne0$, some $y'\in Y$ has $D(f_{\gamma_0})\delta_{y'}\ne0$, and
the coefficient of $\delta_{y'}$ in $\varphi_Y(r)\delta_{\gamma_0^{-1}y'}$
is nonzero. So $\varphi_Y(r)\ne0$. For finite $K\subseteq G$ and large $F$,
entrywise application of $\varphi_Y$ is injective on $K\cup K^{-1}$ and
preserves the products of two elements of $K\cup K^{-1}$. So it maps $K$
into $\GL_{n|Y|}(\F_2)$, injectively and preserving the products that
stay in $K$, and $G$ is LEF.

LEF groups are sofic, and sofic groups are
hyperlinear~\cite[Example~4.5 and Theorem~3.3]{Pestov}. So $L(G)$ embeds
in $\mathcal R^\omega$~\cite[Proposition~7.1]{Ozawa}, where $\mathcal R$ is
the hyperfinite $\mathrm{II}_1$ factor. As $G$ is infinite and simple,
its nontrivial conjugacy classes are infinite, so
$L(G)\mathbin{\bar\otimes}\mathcal R$ is a McDuff factor. It embeds in
$\mathcal R^\omega$ and its unitary group contains $G$, which is Brown's
formulation. Kirchberg proved that a Kazhdan group
with the factorization property is residually
finite~\cite[Theorem~1.1]{Kirchberg}. So $G$ does not have the
factorization property, and $C^*(G)$ does not have the local lifting
property~\cite[p.~527]{Ozawa}.

\subsection*{Simplicity}
Let $1\ne N\trianglelefteq G$ and $1\ne g\in N$, and let $w\ge0$ bound the
word lengths of the group elements occurring in the entries of $g$ and
$g^{-1}$. Call a clopen set $V$ small if $V\cap\gamma V=\varnothing$ for
$0<|\gamma|\le2w$ and every $f\circ\gamma$, with $|\gamma|\le w$ and $f$ a
coefficient of an entry of $g$ or $g^{-1}$, is constant on $V$. The
action is free, so every point has a small clopen neighborhood. By
compactness every clopen set is a finite disjoint union of small ones.

Some $h=e_{ij}(e_V)$ with $V$ small does not commute with $g$.
Otherwise $g$ commutes with $e_{ij}(e_V)$ for every clopen $V$,
by~\eqref{eq:elementary}. Taking $V=X$ gives $g=cI_n$ with
$c=\sum_\gamma c_\gamma u_\gamma\in R$. Then
$e_Vc-ce_V=\sum_\gamma c_\gamma(e_V-e_{\gamma V})u_\gamma$ vanishes for every
clopen $V$. Choosing $V$ to contain $x$ but not $\gamma^{-1}x$, we get
$c_\gamma(x)=0$ for $\gamma\ne e$. So $c\in\LC(X,\F_2)$, and comparing the
coefficients of $u_e$ in $cc^{-1}=1$ gives $c=1$ and $g=1$.

Fix such $h$ and $V$, put $d=n|B_w|$, and for $a,b\in B_w$ put
$\epsilon_{ab}=e_{aV}u_{ab^{-1}}$. Since the sets $aV$, $a\in B_w$, are
disjoint, $\epsilon_{ab}\epsilon_{a'b'}=\delta_{ba'}\epsilon_{ab'}$. Index the
coordinates of $\F_2^d$ by pairs $(p,a)$ with $1\le p\le n$ and
$a\in B_w$. Then $\psi(E_{(p,a),(q,b)})=\epsilon_{ab}E_{pq}$ defines an
injective multiplicative linear map $\psi:M_d(\F_2)\to M_n(R)$, and
$A\mapsto I_n-\psi(I_d)+\psi(A)$ embeds $\GL_d(\F_2)$ in $\GL_n(R)$. Its
image $H$ lies in $G$. Indeed, transvections generate $\GL_d(\F_2)$. The
transvection between $(p,a)$ and $(q,b)$ maps to $e_{pq}(\epsilon_{ab})$
if $p\ne q$. If $p=q$, it is the commutator of the transvections between
$(p,a)$ and $(p',a)$ and between $(p',a)$ and $(p,b)$, where $p'\ne p$.

Put $k=[g,h]\in N\setminus\{1\}$. If $a,b\in B_w$ and $f,f'$ are
coefficients as above, then
\[
  fu_a\,e_V\,f'u_b=f\,e_{aV}\,(f'\circ a^{-1})\,u_{ab}
  \in\{0,\epsilon_{a,b^{-1}}\},
\]
since $f\circ a$ and $f'$ are constant on $V$. The entries of
$ghg^{-1}-I_n=g\,e_VE_{ij}\,g^{-1}$ are sums of such products, and
$h^{-1}=h=I_n+\epsilon_{ee}E_{ij}$. So $k-I_n=(ghg^{-1}-h)h$ and
$k^{-1}-I_n=h(ghg^{-1}-h)$ lie in $\psi(M_d(\F_2))$. As $\psi$ is
injective and multiplicative, $k\in H$. Then $N\cap H$ is a nontrivial
normal subgroup of $H\cong\GL_d(\F_2)=\PSL_d(\F_2)$, which is simple as
$d\ge3$. So $H\subseteq N$, and $e_{pq}(e_V)\in N$ for all $p\ne q$.

Finally, the level $J=\{r\in R:e_{pq}(r)\in N\text{ for all }p\ne q\}$
of $N$~\cite{Stepanov} is a two-sided ideal. It is additive
by~\eqref{eq:elementary}, and for $l\notin\{p,q\}$ we have
$e_{pq}(sr)=[e_{pl}(s),e_{lq}(r)]$ and $e_{pq}(rs)=[e_{pl}(r),e_{lq}(s)]$.
It contains $e_V$, so it contains every $e_{aV}=u_ae_Vu_a^{-1}$. By
minimality and compactness finitely many of these sets cover $X$, so
$1=1-\prod_i(1-e_{a_iV})\in J$. So $J=R$ and $N=G$.\hfill$\square$

\section{Subshifts with finite models}

\begin{proof}[Proof of Proposition~\ref{prop:models}]
(a) Minimality and infiniteness imply that the shift has no periodic
points, so the action of $\Z$ is free. Fix $x\in X$ and $\ell\ge0$. By
minimality every word of $X$ occurs in $x$, and $x_{[0,2\ell)}$ recurs at
arbitrarily large positions. Choose such a position $m$ with all words of
length $2\ell+1$ of $X$ occurring in $x_{[0,m+2\ell)}$. This word begins
and ends with $x_{[0,2\ell)}$, so it has period $m$, and the $m$-periodic
sequence $y$ that agrees with $x$ on $[0,m)$ has the same words of length
$2\ell+1$ as $X$. So the orbit of $y$ is a finite subshift with the
patterns of $X$ on $[-\ell,\ell]$, and so on every subset of $[-\ell,\ell]$.

(b) Let $\gamma\ne e$, and choose $F_\gamma$ as in Section~1, which uses
only freeness. No point of a finite subshift $Y$ with the patterns of $X$
on $F_\gamma\cup F_\gamma\gamma$ is fixed by $\gamma$, so the action of
$\Gamma$ on $Y$ is a homomorphism to a finite group that does not kill
$\gamma$.

(c) Let $\Gamma$ be infinite, finitely generated and residually finite.
Choose normal subgroups $\Gamma=K_0>K_1>\cdots$ of finite index with
$\bigcap_mK_m=\{e\}$ and $[K_m:K_{m+1}]\ge3$, and for $\gamma\ne e$ let
$\ell(\gamma)$ be the $m$ with $\gamma\in K_m\setminus K_{m+1}$. Let
$x\in\{0,1,2\}^\Gamma$ have $x(e)=1$, and on each $K_m\setminus K_{m+1}$ let
$x$ be constant on the cosets of $K_{m+1}$, equal to $1$ on one of them,
to $2$ on another, and to $0$ on the rest. If $h\notin K_n$ and $k\in K_n$,
then $hk$ lies in the coset $hK_{\ell(h)+1}$, so $x(hk)=x(h)$. Let $X$ be
the closure of $\Gamma x$, and call $h\mapsto x(h\sigma)$ on a finite set
$F$ the pattern of $x$ at $\sigma$.

\emph{Finite models.} Let $F\subseteq\Gamma$ be finite, and choose $n$ with
$FF^{-1}\cap K_n=\{e\}$. For $\sigma\in\Gamma$ at most one $h_0\in F$ has
$h_0\sigma\in K_n$. For $k\in K_n$ the patterns of $x$ at $\sigma k$ and at
$\sigma$ agree off $h_0$, and as $k$ varies,
$x(h_0\sigma k)$ takes every value in $x(K_n)$. For $b\in x(K_n)$ let $y_b$
agree with $x$ off $K_n$ and equal $b$ on $K_n$. It is fixed by $K_n$, so
$Y=\bigcup_b\Gamma y_b$ is a finite subshift, and the patterns of $Y$ on
$F$ are exactly those of $x$, which are those of $X$.

\emph{Minimality.} With $F$, $\sigma$, $n$ and $h_0$ as above, the pattern
of $x$ at $\sigma k$ equals the pattern at $\sigma$ for all $k\in K_n$ if
there is no $h_0$, for all $k\in K_{\ell(h_0\sigma)+1}$ if $h_0\sigma\ne e$,
and for all $k$ in the coset of $K_{n+1}$ in $K_n\setminus K_{n+1}$ where
$x=1$ if $h_0\sigma=e$. So every pattern of $x$ recurs along a coset of a
normal subgroup of finite index, $x$ is almost periodic, and $X$ is
minimal.

\emph{Freeness.} Let $\gamma\ne e$, choose $n$ with $\gamma\notin K_n$ and a
finite set $T$ meeting every coset of $K_n$, and let $\sigma\in\Gamma$. Put
$g=\sigma^{-1}\gamma\sigma$ and $m=\ell(g)<n$. Of the two cosets of $K_{m+1}$
in $K_m\setminus K_{m+1}$ where $x=1$ and where $x=2$, one, say
$\kappa K_{m+1}$, differs from $g^{-1}K_{m+1}$. For $\rho\in\kappa K_{m+1}$
the element $\rho g$ lies in another coset of $K_{m+1}$ in
$K_m\setminus K_{m+1}$, so $x(\rho g)\ne x(\rho)$. As $\kappa K_{m+1}\sigma^{-1}$
is a union of cosets of $K_n$, some $h\in T$ has $h\sigma\in\kappa K_{m+1}$,
and then $x(h\sigma)\ne x(h\gamma\sigma)$. So $\sigma x$ and $\gamma\sigma x$
differ on $T$ for every $\sigma$, and by continuity $\gamma z\ne z$ for
every $z\in X$. Finally, $X$ is infinite because $\Gamma$ is infinite and
acts freely.
\end{proof}

The subshifts in (c) are Toeplitz subshifts; see Cortez and
Petite~\cite{CortezPetite} for Toeplitz subshifts over residually finite
groups.

\section{Residually finite subgroups}

\begin{proof}[Proof of Corollary~\ref{cor:rf}]
Let $\Gamma$ be residually finite and generated by a finite set $S$.
Replacing $\Gamma$ by $\Gamma\times\Z$, we may assume that $\Gamma$ is
infinite. Let $\Gamma_m$ be normal subgroups of finite index with trivial
intersection. Left multiplication on $\Gamma/\Gamma_m\times\{1,2\}$ is an
even permutation, so $\Gamma$ embeds in
$P=\prod_m\Sym(\Gamma/\Gamma_m\times\{1,2\})$ with image in the product of
the alternating groups. Every even permutation of a finite set is a
commutator of two permutations~\cite[Theorem~1]{Ore}, so each $s\in S$ is
a commutator $[a_s,b_s]$ in $P$. The subgroup $\Delta$ of $P$ generated by
$S$ and the $a_s,b_s$ is finitely generated, infinite and residually
finite, and $\Gamma\le[\Delta,\Delta]$. By Proposition~\ref{prop:models}(c)
and Theorem~\ref{thm:main}, some $G=\EL_3(\LC(X,\F_2)\rtimes\Delta)$ is an
infinite finitely generated simple Kazhdan LEF group.

For a unit $c$ of a ring of characteristic $2$ we have
$e_{12}(c)e_{21}(c^{-1})e_{12}(c)\,e_{12}(1)e_{21}(1)e_{12}(1)=\diag(c,c^{-1})$,
so for units $a,b$ the matrix
\[
 \diag(a,a^{-1})\diag(b,b^{-1})\diag\bigl((ba)^{-1},ba\bigr)
 =\diag(aba^{-1}b^{-1},1)
\]
lies in $\EL_2$. Since $\Gamma\le[\Delta,\Delta]$, every $u_\gamma$ with
$\gamma\in\Gamma$ is a product of commutators of units, so
$\gamma\mapsto\diag(u_\gamma,1,1)$ is an injective homomorphism
$\Gamma\to G$.
\end{proof}

\section{Word problems}

\begin{corollary}\label{cor:wp}
The word problem of $G_X$ has the Turing degree of the language $L(X)$
of $X$. Every Turing degree occurs, so there are continuum many pairwise
nonisomorphic groups $G_X$.
\end{corollary}

\begin{proof}
Multiplying out a word in the generators gives a matrix with entries
$\sum_jf_ju^j$, each $f_j$ a table on words, using $uf=(f\circ T^{-1})u$.
The word is trivial if and only if the tables of its difference from
$I_3$ vanish on $L(X)$, so $L(X)$ computes the word problem. Conversely,
by~\eqref{eq:elementary} a word for $e_{12}(\prod_{t<n}u^{-t}e_{v_t}u^t)$
is computable from $v=v_0\cdots v_{n-1}$, and it is trivial if and only
if $v\notin L(X)$. For derived topological full groups, Grigorchuk and
Medynets proved that the word problem is decidable if and only if
$L(X)$ is recursive~\cite[Theorem~1.1(3)]{GMpres}.

For irrational $\alpha\in(0,1)$ let $X_\alpha$ be the infinite minimal
Sturmian subshift, the closure of the codings $c(\theta)$,
$\theta\in[0,1)$, with $c(\theta)_t=1$ if and only if
$\theta+t\alpha\bmod1\ge1-\alpha$~\cite{MorseHedlund}. Its words of
length $n$ are the constant values of $c(\theta)_{[0,n)}$ on the arcs
between the points $-j\alpha\bmod1$, $0\le j\le n$, so $\alpha$ computes
$L(X_\alpha)$. The word $c(\theta)_{[0,n)}$ has
$\lfloor\theta+n\alpha\rfloor$ ones, within $1$ of $n\alpha$, so
$L(X_\alpha)$ computes $\alpha$. Every set $Q\subseteq\mathbb N$ has the
degree of the irrational number $[0;1+\chi_Q(0),1+\chi_Q(1),\dots]$, and
the degree of the word problem is an isomorphism invariant.
\end{proof}

\section{Questions}

A finitely presented LEF group is residually finite~\cite{VershikGordon},
and an infinite simple group is not, so the groups of
Theorem~\ref{thm:main} are not finitely presented.

\begin{question}
Is there a finitely presented infinite simple group with
property~\textup{(T)} that is sofic, or at least hyperlinear?
\end{question}

A positive answer would also answer Open problem~6.1 of Alekseev and
Thom~\cite{AlekseevThom}, which asks for a finitely presented sofic group
with property~\textup{(T)} that is not residually finite.

\begin{question}
Is every finitely generated LEF group a subgroup of an infinite simple
Kazhdan LEF group?
\end{question}

If $(X,T)$ is conjugate to $(Y,S)$ or to $(Y,S^{-1})$, then $G_X\cong G_Y$.

\begin{question}
Does $G_X\cong G_Y$ imply that $X$ and $Y$ are flip conjugate, or at least
strongly orbit equivalent~\cite{GPS}?
\end{question}

\subsection*{Origin and authorship}
Claude (Anthropic) found the construction and original proofs under
the author's direction on September~12, 2026, and wrote the original
manuscript. Codex (OpenAI) shortened an earlier version, and Claude
revised this one.
The author is responsible for the final manuscript. Proof records are
in the \href{https://github.com/SauersML/group-approximation}{project repository}.

\begin{thebibliography}{99}

\bibitem{AlekseevThom}
V.~Alekseev and A.~Thom,
\emph{Centralizers of sofic approximations of Kazhdan groups},
\href{https://arxiv.org/abs/2608.05362}{arXiv:2608.05362} (2026).

\bibitem{Brown}
N.~P. Brown, \emph{Tracial invariants, classification and
$\mathrm{II}_1$ factor representations of Popa algebras},
\href{https://arxiv.org/abs/math/0111286}{arXiv:math/0111286} (2001).

\bibitem{CortezPetite}
M.~I. Cortez and S.~Petite,
\emph{$G$-odometers and their almost one-to-one extensions},
J. Lond. Math. Soc. (2) \textbf{78} (2008), 1--20.
\doi{10.1112/jlms/jdn002}.

\bibitem{EJZ}
M.~Ershov and A.~Jaikin-Zapirain,
\emph{Property~\textup{(T)} for noncommutative universal lattices},
Invent. Math. \textbf{179} (2010), 303--347.
\doi{10.1007/s00222-009-0218-2}.

\bibitem{GPS}
T.~Giordano, I.~F. Putnam, and C.~F. Skau,
\emph{Topological orbit equivalence and $C^*$-crossed products},
J. Reine Angew. Math. \textbf{469} (1995), 51--112.
\doi{10.1515/crll.1995.469.51}.

\bibitem{GM}
R.~Grigorchuk and K.~Medynets,
\emph{On algebraic properties of topological full groups},
Sb. Math. \textbf{205} (2014), 843--861.
\doi{10.1070/SM2014v205n06ABEH004400}.

\bibitem{GMpres}
R.~Grigorchuk and K.~Medynets,
\emph{Presentations of topological full groups by generators and
relations}, J. Algebra \textbf{500} (2018), 46--68.
\doi{10.1016/j.jalgebra.2016.10.027}.

\bibitem{KerrNowak}
D.~Kerr and P.~W. Nowak,
\emph{Residually finite actions and crossed products},
Ergodic Theory Dynam. Systems \textbf{32} (2012), 1585--1614.
\doi{10.1017/S0143385711000575}.

\bibitem{KionkeSchesler}
S.~Kionke and E.~Schesler,
\emph{From telescopes to frames and simple groups},
J. Comb. Algebra (2024). \doi{10.4171/JCA/103}.

\bibitem{Kirchberg}
E.~Kirchberg, \emph{Discrete groups with Kazhdan's property~$T$ and
factorization property are residually finite},
Math. Ann. \textbf{299} (1994), 551--563.
\doi{10.1007/BF01459798}.

\bibitem{Matui}
H.~Matui, \emph{Some remarks on topological full groups of Cantor
minimal systems}, Internat. J. Math. \textbf{17} (2006), 231--251.
\doi{10.1142/S0129167X06003448}.

\bibitem{MorseHedlund}
M.~Morse and G.~A. Hedlund,
\emph{Symbolic dynamics II. Sturmian trajectories},
Amer. J. Math. \textbf{62} (1940), 1--42.
\doi{10.2307/2371431}.

\bibitem{Ore}
O.~Ore, \emph{Some remarks on commutators},
Proc. Amer. Math. Soc. \textbf{2} (1951), 307--314.
\doi{10.1090/S0002-9939-1951-0040298-4}.

\bibitem{Ozawa}
N.~Ozawa, \emph{About the QWEP conjecture},
Internat. J. Math. \textbf{15} (2004), 501--530.
\doi{10.1142/S0129167X04002417}.

\bibitem{Pestov}
V.~G. Pestov, \emph{Hyperlinear and sofic groups: a brief guide},
Bull. Symbolic Logic \textbf{14} (2008), 449--480.
\doi{10.2178/bsl/1231081461}.

\bibitem{Stepanov}
A.~V. Stepanov, \emph{On the normal structure of the general linear
group over a ring}, Zap. Nauchn. Sem. POMI \textbf{236} (1997), 166--182;
English transl., J. Math. Sci. \textbf{95} (1999), 2146--2155.
\doi{10.1007/BF02169976}.

\bibitem{Thom}
A.~Thom, \emph{Examples of hyperlinear groups without factorization
property}, Groups Geom. Dyn. \textbf{4} (2010), 195--208.
\doi{10.4171/GGD/80}.

\bibitem{VershikGordon}
A.~M. Vershik and E.~I. Gordon,
\emph{Groups that are locally embeddable in the class of finite groups},
Algebra i Analiz \textbf{9} (1997), no.~1, 71--97; English transl.,
St. Petersburg Math. J. \textbf{9} (1998), no.~1, 49--67.

\end{thebibliography}
\end{document}
